% BHU PYQ -- Manually transcribed & error-corrected
\documentclass[11pt,a4paper]{article}

\usepackage[a4paper,top=2.5cm,bottom=2.5cm,left=2.8cm,right=2.8cm,headheight=14pt]{geometry}
\usepackage{amsmath,amssymb}
\usepackage[T1]{fontenc}
\usepackage[utf8]{inputenc}
\usepackage{lmodern}
\usepackage{microtype}
\usepackage{fancyhdr}
\usepackage{enumitem}
\usepackage{hyperref}
\usepackage{lastpage}
\usepackage{parskip}

\hypersetup{colorlinks=true,urlcolor=black,linkcolor=black,
  pdftitle={BPT-602: Solid State Physics -- BHU PYQ},pdfauthor={Banaras Hindu University}}

\pagestyle{fancy}\fancyhf{}
\renewcommand{\headrulewidth}{0.4pt}\renewcommand{\footrulewidth}{0.4pt}
\fancyhead[L]{\small   \textit{Banaras Hindu University}}
\fancyhead[R]{\small   \textit{BPT-602: Solid State Physics}}
\fancyfoot[C]{\small Page  \thepage\ of \pageref{LastPage}}


\newlist{parts}{enumerate}{2}
\setlist[parts,1]{label=(\alph*),leftmargin=*,itemsep=5pt,topsep=3pt}
\setlist[parts,2]{label=(\roman*),leftmargin=*,itemsep=3pt,topsep=2pt}

\newcommand{\pts}[1]{\hfill   \textit{[#1]}}


\begin{document}


\begin{center}
  {\large\bfseries BANARAS HINDU UNIVERSITY}\\[2pt]
  {
\normalsize B.Sc. (Hons.) Semester VI Examination, 2013-14}\\[6pt]
  \rule{0.8\linewidth}{0.5pt}\\[6pt]
  {\large\bfseries Paper No. BPT-602: Solid State Physics}\\[4pt]
  {
\normalsize Time: 3 Hours\hfill Maximum Marks: 70}
\end{center}
\rule{\linewidth}{0.4pt}
\smallskip



\noindent   \textit{Instructions: Answer all questions. Symbols have their usual meanings.}
\bigskip




\noindent   \textbf{Question 1.} Answer any   \textbf{five} of the following questions: \pts{5 $ \times$ 4 = 20}

\begin{parts}
  \item With the help of diagrams, show the crystal structure of NaCl and diamond. What is the lattice type of these structures?
  \item For metals, establish the Wiedemann-Franz law.
  \item Show the variation of energy, velocity, and effective mass of an electron as a function of wave vector $k$ in the first Brillouin zone.
  \item What are the characteristic differences between the optical and acoustical modes of lattice vibrations?
  \item What do you understand by the Lennard-Jones potential?
  \item Explain the meaning of spin waves and magnons.
\end{parts}

\medskip\rule{\linewidth}{0.2pt}




\noindent   \textbf{Question 2.}

\begin{parts}
  \item What is the meaning of a space-lattice? How many different types of space-lattices are possible? \pts{6}
  \item Prove that the close packing of atoms in the hexagonal close-packed (HCP) structure demands an axial ratio $c/a = \sqrt{8/3} \approx 1.633$. \pts{6}
\end{parts}

\medskip\rule{\linewidth}{0.2pt}




\noindent   \textbf{Question 3.}

\begin{parts}
  \item Using the Einstein model, derive the expression for the specific heat of solids at high and low temperatures. What are the assumptions of the Debye model of specific heat of solids? How does it differ from the Einstein model? \pts{8}
  \item Considering the nearest neighbour interaction, establish the dispersion relation for a monoatomic linear chain. \pts{5}
\end{parts}
\hfill   \textbf{OR}

\begin{parts}
  \item Derive the expression for the energy of conduction electrons in a metal considering its potential energy to be zero. \pts{8}
  \item On the basis of the free electron gas model, calculate the expression for electrical conductivity. \pts{5}
\end{parts}

\medskip\rule{\linewidth}{0.2pt}




\noindent   \textbf{Question 4.}

\begin{parts}
  \item Discuss the quantum theory of paramagnetism and establish the Curie-Brillouin law for the magnetization and hence the Curie law of susceptibility. \pts{8}
  \item What is the Weiss field for ferromagnetism? Establish the Curie-Weiss law for ferromagnetism. \pts{5}
\end{parts}
\hfill   \textbf{OR}

\begin{parts}
  \item What do you understand by reciprocal lattice? \pts{2}
  \item Show that the reciprocal vector $\vec{G}_{hkl}$ is perpendicular to the $(hkl)$ plane and $|\vec{G}_{hkl}| = \frac{2\pi}{d_{hkl}}$. \pts{6}
  \item What is the Ewald construction? \pts{5}
\end{parts}

\medskip\rule{\linewidth}{0.2pt}




\noindent   \textbf{Question 5.}

\begin{parts}
  \item Derive the expressions for the atomic scattering factor and crystal structure factor. \pts{8}
  \item Calculate the structure factor for a body-centered cubic (BCC) structure and show that its value is non-zero only when $h+k+l$ is an even number, where $h$, $k$, and $l$ are the Miller indices of the plane. \pts{4}
\end{parts}

\vfill
\rule{\linewidth}{0.4pt}

\begin{center}\small
  \href{https://bhu-exam.vercel.app/}{Click here to visit our website}\\[4pt]
  \href{https://me-aryan.vercel.app/}{Click here to visit our contributor}\\[4pt]
  \href{https://quashan-me.vercel.app/}{Click here to visit our contributor}
\end{center}

\end{document}
